% !TeX root = ../main.tex
\section{Изменение количества бит цифрового генератора}
Проанализируем изменение параметра бит цифрового гетеродина $ n_{g} $, при условии что тип бит это с фиксированной запятой от 2 до 11 для моего варианта.

\begin{figure}[H]
    \centering
    \begin{minipage}[t]{0.48\textwidth}
        \vspace{0pt}
        \centering
        \includegraphics[width=\linewidth]{3_spec_ng2.png}
        \caption{Спектр на выходе DEM ($n_g = 2$)}
        \label{fig:3_ng2}
    \end{minipage}
    \hfill
    \begin{minipage}[t]{0.48\textwidth}
        \vspace{0pt}
        \centering
        \includegraphics[width=\linewidth]{3_spec_ng4.png}
        \caption{Спектр на выходе DEM ($n_g = 4$)}
        \label{fig:3_ng4}
    \end{minipage}
\end{figure}

Анализируя графики представленные выше, можно заметить, что количество побочных гармоник снизилось. Это случилось из-за того что при большей разрядности цифрового гетеродина ошибки округления становятся более редкими, точность перемножения увеличивается, вследствие чего гармоник становится меньше. Проверим данную гипотезу при других разрядностях (6,9,11).

\begin{figure}[H]
    \centering
    \begin{minipage}[t]{0.48\textwidth}
        \vspace{0pt}
        \centering
        \includegraphics[width=\linewidth]{3_spec_ng6.png}
        \caption{Спектр на выходе DEM ($n_g = 6$)}
        \label{fig:3_ng6}
    \end{minipage}
    \hfill
    \begin{minipage}[t]{0.48\textwidth}
        \vspace{0pt}
        \centering
        \includegraphics[width=\linewidth]{3_spec_ng9.png}
        \caption{Спектр на выходе DEM ($n_g = 9$)}
        \label{fig:3_ng9}
    \end{minipage}
\end{figure}

\begin{figure}[H]
    \centering
    \includegraphics[width=0.8\textwidth]{3_spec_ng11.png}
    \caption{Спектр на выходе DEM ($n_g = 11$)}
\end{figure}

Сравним результаты при 2 и 11 количества бит. Количество бит напрямую связывает количество паразитных составляющих, которые появляются после перемножения. Гипотеза об округлении верна.